% !TeX root = ../../Book1.tex
% 纲要标题：### 20.1 循环扩张
% 纲要位置：计划纲要.md，第 676 行。

\section{循环扩张}
\label{b1:sec:2001}

循环扩张的全部自同构都由一个自同构反复复合得到。因此，迹与范数分别成为沿一个有限轨道求和、求积的运算。我们要问的是：这些运算的核能否用同一个生成自同构直接描述？先从容易验证的一个方向入手，再在下一节证明其逆命题。

% 纲要内容（仅作写作计划，尚非教材正文）：
%
% * 循环 Galois 群与生成自同构
% * 范数为一的元素
% * 加法与乘法形式的区别
%

\subsection{循环 Galois 群与生成自同构}
\label{b1:subsec:2001:01}

\begin{definition}\label{b1:def:cyclic-extension}
设 \(L/K\) 为有限 Galois 扩张。如果保持 \(K\) 中每个元素不变的全部 \(L\) 自同构组成的群 \(\operatorname{Gal}(L/K)\) 是循环群，则称 \(L/K\) 为\term[xunhuan-kuozhang]{循环扩张}[cyclic extension]。记 \([L:K]=n\)，并选取 \(\operatorname{Gal}(L/K)\) 的一个生成元 \(\sigma\)，于是
\[
G=\operatorname{Gal}(L/K)=\{1,\sigma,\ldots,\sigma^{n-1}\},
\qquad \sigma^n=1.
\]
\end{definition}

这里必须同时要求正规与可分。一个扩张的自同构群偶然是循环群，并不能保证它是循环扩张。例如 \(\mathbb Q(\sqrt[3]{2})/\mathbb Q\) 的自同构群只有恒等映射，但它不是正规扩张。相反，在特征不为 \(2\) 的域上，非平凡二次扩张 \(K(\sqrt d)/K\) 的两个自同构为恒等映射及 \(\sqrt d\mapsto-\sqrt d\)，故为循环扩张。

\ProseParagraphBreak
生成元通常不是唯一的。\(\sigma^r\) 生成同一个 \(n\) 阶群当且仅当 \(\gcd(r,n)=1\)。由\cref{b1:thm:finite-galois-correspondence}，每个 \(n\) 的正因子 \(d\) 对应唯一的 \(d\) 次中间扩张；具体地，\(L^{\langle\sigma^d\rangle}\) 在 \(K\) 上的次数为 \(d\)。这个次数公式使用的是子群指数，而不是子群阶。

\subsection{范数为一的元素}
\label{b1:subsec:2001:02}

\cref{b1:thm:trace-norm-embeddings}在当前情形给出
\begin{equation}\label{b1:eq:cyclic-trace-norm}
\operatorname{Tr}_{L/K}(x)=\sum_{i=0}^{n-1}\sigma^i(x),
\qquad
N_{L/K}(x)=\prod_{i=0}^{n-1}\sigma^i(x).
\end{equation}
特别地，若 \(b\in L^\times\)，则
\[
N_{L/K}\Bigl(\frac{b}{\sigma(b)}\Bigr)
=\prod_{i=0}^{n-1}\frac{\sigma^i(b)}{\sigma^{i+1}(b)}=1.
\]
分子、分母逐项相消，末项使用 \(\sigma^n(b)=b\)。因此每个这样的比值都在范数的核中。

\ProseParagraphBreak
以 \(\mathbb C/\mathbb R\) 为例，生成自同构是复共轭，范数为 \(z\bar z\)。非零数 \(b/\bar b\) 总在单位圆上；反过来，单位圆上的每个数能否如此表示，就是乘法形式要回答的问题。这个问题只要求找出一个 \(b\)，不要求唯一：用任意非零实数乘 \(b\)，比值并不改变。

\subsection{加法形式与乘法形式}
\label{b1:subsec:2001:03}

对加法作同样计算，可得
\[
\operatorname{Tr}_{L/K}\Bigl(b-\sigma(b)\Bigr)
=\sum_{i=0}^{n-1}\Bigl(\sigma^i(b)-\sigma^{i+1}(b)\Bigr)=0.
\]
两种计算分别给出
\[
\Bigl\{\,b/\sigma(b)\mid b\in L^\times\,\Bigr\}\subseteq\ker N_{L/K},
\qquad
\{\,b-\sigma(b)\mid b\in L\,\}\subseteq\ker\operatorname{Tr}_{L/K}.
\]
乘法形式的目标群是 \(L^\times\)，其中零没有意义；加法形式的目标是 \(L\) 的加法群。加法形式也不能在特征整除 \(n\) 时靠除以 \(n\) 来证明。下一节的论证对这些特征仍然成立，因为真正使用的是扩张的可分性。
